% Created 2026-06-27 Sat 15:17
% Intended LaTeX compiler: pdflatex
\documentclass[11pt,letterpaper]{article}
\usepackage[utf8]{inputenc}
\usepackage[T1]{fontenc}
\usepackage{graphicx}
\usepackage{longtable}
\usepackage{wrapfig}
\usepackage{rotating}
\usepackage[normalem]{ulem}
\usepackage{amsmath}
\usepackage{amssymb}
\usepackage{capt-of}
\usepackage{hyperref}
\usepackage{amsmath}
\usepackage{amssymb}
\usepackage{geometry}
\geometry{left=1.2in, right=1.2in, top=1.0in, bottom=1.0in}
\usepackage{fancyhdr}
\usepackage{booktabs}
\usepackage{array}
\usepackage{parskip}
\usepackage{microtype}
\usepackage{tcolorbox}
\tcbuselibrary{skins,breakable}
\newenvironment{plainbox}{\begin{tcolorbox}[colback=blue!3!white,colframe=blue!40!white,title={\small\textit{In plain language}},breakable]}{\end{tcolorbox}}
\newenvironment{openbox}{\begin{tcolorbox}[colback=orange!5!white,colframe=orange!60!black,title={\small\textit{Open calculation}},breakable]}{\end{tcolorbox}}
\newenvironment{defbox}{\begin{tcolorbox}[colback=green!3!white,colframe=green!45!black,title={\small\textit{Foundational definition}},breakable]}{\end{tcolorbox}}
\PassOptionsToPackage{hidelinks}{hyperref}
\setlength{\parindent}{0pt}
\setlength{\parskip}{6pt}
\fancypagestyle{plain}{\fancyhf{}\fancyfoot[C]{\thepage}\renewcommand{\headrulewidth}{0pt}}
\fancypagestyle{body}{%
\fancyhf{}%
\fancyhead[C]{\small Document 2 of 5 $\cdot$ DOI: 10.5281/zenodo.18927154 $\cdot$ CC BY 4.0}%
\fancyfoot[C]{\thepage}%
\renewcommand{\headrulewidth}{0.4pt}}
\author{Bradley K. DePuy, Independent Researcher}
\date{2026}
\title{Negotiation Modes, Particle Structure, and the Derivation of the Proton Mass}
\hypersetup{
 pdfauthor={Bradley K. DePuy, Independent Researcher},
 pdftitle={Negotiation Modes, Particle Structure, and the Derivation of the Proton Mass},
 pdfkeywords={},
 pdfsubject={},
 pdfcreator={Emacs 29.3 (Org mode 9.6.15)}, 
 pdflang={English}}
\begin{document}

\pagestyle{plain}
\begin{titlepage}
\vspace*{3cm}
\begin{center}
  {\large\bfseries THE MEDIUM --- AN ATTEMPTED DESCRIPTION}\\[1em]
  {\LARGE\bfseries Negotiation Modes, Particle Structure,}\\[0.4em]
  {\LARGE\bfseries and the Derivation of the Proton Mass}\\[0.4em]
  {\large Structure from Substrate Primitives; Scale from Observation}\\[0.3em]
  {\normalsize Version 3 --- with Foundational Iota Scale Definition}\\[2em]
  {\normalsize Bradley K. DePuy, Independent Researcher}\\[0.4em]
  {\normalsize ORCID: \texttt{0009-0001-0328-2299}}\\[0.4em]
  {\normalsize Zenodo DOI: 10.5281/zenodo.18927154}\\[0.4em]
  {\normalsize github.com/bkdepuy/The-Medium}\\[0.4em]
  {\normalsize 2026}
\end{center}
\end{titlepage}
\pagestyle{body}

\section{Abstract}
\label{sec:orgf7efeca}

This paper develops a connected chain of derivation within The Medium, from fundamental primitives through particle structure to the proton mass and the negotiation energy quantum, with open calculations explicitly identified at each step. A foundational definition is established at the outset: iota size \(1\) corresponds to the Planck length. This makes the Planck scale definitional rather than empirical, reduces the external inputs to exactly two --- the iota scale and the Hubble scale --- and resolves the previously noted circularity in the derivation of the Planck constant \(\hbar\).

Beginning from this foundation and the primitive entities --- iotas (\(\iota\)) and their boundary negotiation dynamics --- we establish that quantization, confinement, and a three-fold mass ratio follow necessarily from the topology of closed negotiation graphs. The minimum closed negotiation graph is a triangle; its Laplacian eigenvalue \(\lambda = 3\) predicts that any stable configuration of this topology has energy \(E = 3\varepsilon\), where \(\varepsilon\) is the fundamental negotiation energy quantum. Matching to the measured proton mass gives \(\varepsilon = 312.8\) MeV --- this is a definition of \(\varepsilon\) from observation, not a derivation of the proton mass from the substrate. Independently, the geometric mean of the two scale energies gives \(\varepsilon_\text{est} \approx 152\) MeV, a factor of \({\sim}2\) below the observed value. Whether this factor is the gluon (edge) contribution or a different correction is an open question. Color charge as an emergent label and spin as relational closure are structural consequences of the triangle topology; their precise connection to the QCD quantum numbers is an open derivation. Testable predictions include the hadron mass ratios and a slow cosmological evolution of the proton mass.

\begin{defbox}
\textbf{Foundational definition:} Iota size $1 =$ Planck length $= 1.616 \times 10^{-35}$ m.

Two external inputs: the iota scale and the Hubble scale $R_H = c/H_0$.

In Planck units: $G = 1$, $c = 1$, $\hbar = 1/(2\pi)$. $\hbar$, $G$, $m_\text{proton}$, $H_0$, and $\Lambda$ are all derived from these two inputs plus the Barbero-Immirzi correction $\gamma \approx 0.2375$.
\end{defbox}

\section{Part I: Accelerators and the Nature of Structure}
\label{sec:org486888c}

\begin{plainbox}
Particle accelerators probe structure at the smallest scales. At modest energies, collisions reveal stable negotiation patterns --- configurations satisfying the closure condition $L \cdot N = 0$. At extreme energies, collisions force configurations the substrate cannot sustain: they decay immediately because they are not stable attractors in the negotiation geometry. The Standard Model does not distinguish these two cases. The Medium does.
\end{plainbox}

\subsection{Accelerators as Negotiation Probes}
\label{sec:orgafa23be}

In quantum mechanics every particle carries a de Broglie wavelength, and resolution is limited by that wavelength. Higher energy produces shorter wavelength, enabling finer structural probing. In The Medium, probing at modest energies reads stable negotiation patterns. Probing at extreme energies forces configurations the substrate cannot sustain, which decay because they are not stable attractors in the negotiation geometry. The Standard Model is agnostic about this distinction. The Medium is not: stable particles are closed negotiation graphs; transient phenomena are off-closure configurations that the path integral does not sustain.

\subsection{The Iota Scale and Particle Size}
\label{sec:orgf370d43}

With iota size \(1 =\) Planck length, the size of a particle is determined by how many iota scales its negotiation pattern spans. A proton's charge radius is \({\sim}0.85 \times 10^{-15}\) m, which is:

\begin{equation}
r_\text{proton} / l_\text{Planck} = 0.85 \times 10^{-15} / 1.616 \times 10^{-35} \approx 5.3 \times 10^{20}\ \text{iota units}
\end{equation}

The proton is not three iotas in physical size. It is three iotas in negotiation topology --- the minimum closed graph. Its physical size emerges from the negotiation wavelength at energy \(\varepsilon\), which spans \({\sim}10^{20}\) iota units. The topology is minimal; the physical extent is large.

\section{Part II: Negotiation Modes --- The Formal Structure}
\label{sec:orge86758f}

\begin{plainbox}
A negotiation between two iotas establishes a boundary. For a negotiation to be complete --- for every boundary to be resolved with no loose ends --- you need at least three iotas in a closed loop. That is a triangle. The negotiation field dynamics are governed by \(\partial^2 N/\partial t^2 = -LN\), whose solutions are normal modes \(N_k(t) \propto v_k e^{i\omega_k t}\) with frequencies \(\omega_k = \sqrt{\lambda_k}\). The \(\lambda = 0\) mode (null space of \(L\)) is the uniform ground state: all boundary tensions equal, no differentiation, no particle. Stable particles are the persistent \textit{non-trivial} normal modes --- the eigenvectors at \(\lambda > 0\) that oscillate stably without decaying.
\end{plainbox}

\subsection{The Relational Graph}
\label{sec:orgf283049}

Represent an iota configuration as a graph \(G\) where each iota is a node, each active negotiation is an edge, and each edge carries a negotiation state \(N_{ij}\). In Planck units, \(N_{ij}\) is dimensionless --- the boundary resolution expressed as a ratio in iota-scale units. A negotiation mode is a pattern of edge states satisfying self-consistency, stability, and closure.

\subsection{The Closure Condition}
\label{sec:org82c3be2}

For a node \(\iota_i\) with neighbors \(\iota_j\), self-consistency requires that every iota's net boundary tension sums to zero:

\begin{equation}
\sum_j N_{ij} = 0 \qquad \Longleftrightarrow \qquad L \cdot N = 0
\end{equation}

The null space of \(L\) (solutions to \(L \cdot N = 0\)) is the uniform mode --- the vacuum. Stable particles correspond to the non-zero eigenspace: \(L \cdot N = \lambda N\) with \(\lambda > 0\), which gives oscillatory normal modes. In Planck units these eigenvalues are dimensionless ratios; their product with the energy quantum \(\varepsilon\) gives particle masses.

\subsection{Quantization from Closure}
\label{sec:org98c1a02}

For a closed negotiation loop of \(n\) iotas, the self-consistency condition requires:

\begin{equation}
n \cdot \phi = 2\pi k \qquad \Longrightarrow \qquad \phi = \frac{2\pi k}{n},\quad k \in \mathbb{Z}
\end{equation}

Only integer mode numbers are permitted. Quantization is a topological necessity, not an added postulate. In Planck units, the phase \(\phi\) is measured in units of the minimum negotiation cycle --- one Planck time.

\section{Part III: The Triangle --- Minimum Closed Structure}
\label{sec:org26a9eff}

\begin{plainbox}
The minimum closed negotiation requires three iotas --- two cannot close, four carry higher energy. The triangle topology is the unique ground-state solution. What the mathematics produces from this topology: \textbf{a three-fold mass ratio} (eigenvalue \(\lambda = 3\)), \textbf{a doubly degenerate mode} (two independent eigenvectors at the same energy), \textbf{three distinguishable node positions} (relational structure of the closed loop), and \textbf{structural confinement} (a node cannot exist outside the graph that defines it). These are mathematical facts about the triangle Laplacian. Their identification with color charge, spin-$\frac{1}{2}$, and quark structure is a structural correspondence with QCD, not yet a derivation from substrate dynamics alone.
\end{plainbox}

\subsection{Why Three}
\label{sec:orgc560a74}

Two iotas negotiating produce one edge but leave each iota with an unresolved side. The minimum closed structure requires three. In Planck units, the energy cost of closing the negotiation is minimized at three nodes. Fewer nodes cannot achieve closure. More nodes carry higher energy. Three is the ground state of the closure condition.

\subsection{The Triangle Laplacian in Planck Units}
\label{sec:org0107ad9}

With iota size \(1 =\) Planck length, the adjacency matrix and Laplacian are dimensionless:

\begin{equation}
A = \begin{pmatrix} 0 & 1 & 1 \\ 1 & 0 & 1 \\ 1 & 1 & 0 \end{pmatrix}, \qquad
L = 2I - A = \begin{pmatrix} 2 & -1 & -1 \\ -1 & 2 & -1 \\ -1 & -1 & 2 \end{pmatrix}
\end{equation}

Eigenvalue spectrum: \(\{\lambda_0 = 0,\ \lambda_1 = 3,\ \lambda_2 = 3\}\). These are pure numbers --- ratios in iota units.

\subsection{Mode Spectrum and Physical Properties}
\label{sec:org9df911a}

\begin{itemize}
\item \(\lambda_0 = 0\) --- ground mode; uniform state; no differentiation; no particle.
\item \(\lambda_1 = \lambda_2 = 3\) (doubly degenerate) --- the two independent non-trivial eigenvectors of the triangle. These are the stable oscillatory modes. They are structurally identified with the proton's two spin states. Whether they transform as a spin-\(\tfrac{1}{2}\) doublet under SU(2) requires deriving how the rotation group acts on the triangle's eigenvectors and showing these modes carry the fundamental representation. That derivation has not been done.
\end{itemize}

The three distinguishable node positions within the closed loop are color charge. The impossibility of removing a node from the graph that defines it is confinement. Neither is postulated.

\subsection{Confinement, Quarks, and Gluons}
\label{sec:orgb1e5f03}

The three nodes of the triangle are structurally analogous to quarks --- aspects of the relational structure, not independent objects. The three edges are structurally analogous to gluons. The mathematics requires both: a graph has both nodes and edges. What the math does not yet provide is the energy carried by each edge independently of the nodes. QCD reports that gluons carry approximately 50$\backslash$% of the proton's mass; the geometric mean formula gives \(\varepsilon_\text{est} \approx \varepsilon/2\). These agree numerically. Whether the factor of 2 is the edge (gluon) energy or a dispersion correction is undetermined --- both are candidates, neither is established.

\section{Part IV: Deriving the Negotiation Energy Quantum \(\varepsilon\)}
\label{sec:orgfe3e43b}

\begin{plainbox}
With iota size $1 =$ Planck length, the derivation of $\varepsilon$ has a clean physical meaning. The proton's negotiation quantum is the geometric mean of two iota-scale energies (one per iota participant in the binary negotiation) and one Hubble-scale energy (the universal boundary condition). The cube root comes from having three energy inputs. The result lands within 10.3\% of the measured value using only the iota scale definition and the size of the observable universe.
\end{plainbox}

\subsection{The Energy Equation}
\label{sec:org60ff6bf}

For a complete graph on \(n\) nodes with mode number \(k\), the Hamiltonian eigenvalue equation gives:

\begin{equation}
E(n,k) = \varepsilon \cdot k \cdot n
\end{equation}

The topology determines the ratio \(n\): the triangle Laplacian has non-trivial eigenvalue \(\lambda = 3\), so the triangle mode carries three times the fundamental quantum. What the topology does not determine is the value of \(\varepsilon\) itself. \(\varepsilon\) is the energy scale of one negotiation cycle --- it requires an independent determination. Matching to the measured proton mass:

\begin{equation}
\varepsilon \equiv m_\text{proton} c^2 / 3 = 938.3\ \text{MeV} / 3 = 312.8\ \text{MeV}
\end{equation}

This is a definition, not a derivation. What this description has established: if the proton is the minimum closed negotiation graph, its mass must be exactly three times the fundamental quantum. The task remaining is to derive \(\varepsilon\) from the substrate scale inputs alone.

\subsection{Estimating \(\varepsilon\) from the Two Scale Inputs}
\label{sec:org5998f99}

The Hubble floor energy in Planck units:

\begin{equation}
E_\text{Hubble} = E_\text{Planck} \cdot \frac{l_\text{Planck}}{R_H} = 9.19 \times 10^{-39}\ \text{MeV}
\end{equation}

The geometric mean of the two scale energies with a 2:1 weighting:

\begin{equation}
\varepsilon_\text{est} = \bigl(E_\text{Planck}^2 \cdot E_\text{Hubble}\bigr)^{1/3}
\end{equation}

This formula is an empirical observation: the quantity \((E_P^2 E_H)^{1/3}\) lands near the observed energy scale of nuclear physics. The 2:1 exponent ratio is not derived from substrate dynamics --- it is the value that makes the formula work. A physical interpretation (two local iota participants, one global boundary condition) has been proposed, but that interpretation was constructed to rationalize the exponent after the fact. The derivation of the weighting from first principles is open.

\subsection{Numerical Result}
\label{sec:org934fd92}

\begin{align}
\varepsilon_\text{predicted} &= \bigl((1.956\times10^{22})^2 \cdot (9.19\times10^{-39})\bigr)^{1/3} = 152.6\ \text{MeV} \notag\\
\varepsilon_\text{corrected} &= (2\pi)^{1/3} \cdot 152.6 = 281.5\ \text{MeV}
\end{align}

Observed \(\varepsilon\) from proton mass: 312.8 MeV. Estimate from scale inputs: 281.5 MeV. Ratio: 312.8/281.5 = 1.11. The \(\sim\)10$\backslash$% residual gap has not been derived. The Barbero-Immirzi parameter \(\gamma \approx 0.2375\) from loop quantum gravity has a magnitude that would close this gap if applied as a dispersion correction (\(\omega = \lambda\) rather than \(\omega = \sqrt{\lambda}\)). Whether \(\gamma\) is the right correction, and why, requires a derivation from the negotiation graph spectrum that does not yet exist.

\subsection{Hadron Mass Predictions}
\label{sec:org601ad51}

General formula \(E(n,k) = \varepsilon \cdot k \cdot n\) for complete graphs:

\begin{center}
\begin{tabular}{lllll}
\toprule
\textbf{Structure} & \(n\) & \textbf{Predicted} & \textbf{Observed} & \textbf{Gap}\\[0pt]
\midrule
Proton & 3 & 938.3 MeV & 938.3 MeV & Exact (by construction)\\[0pt]
Neutron & 3+\(\delta\) & 939.6 MeV & 939.6 MeV & \(\delta = 0.00413\)\\[0pt]
Rho meson & 2 & 625.6 MeV & 775 MeV & 19\% (gluon correction pending)\\[0pt]
Tetraquark & 4 & 1251 MeV & \(\sim\)1200--1400 MeV & Edge corrections needed\\[0pt]
Pion & --- & Not produced & 135 MeV & Not explained by this formula\\[0pt]
\bottomrule
\end{tabular}
\end{center}

The pion is the lightest hadron and the most important test case. It is a pseudo-Goldstone boson of chiral symmetry breaking and cannot be accommodated by the complete-graph formula \(E(n,k) = \varepsilon \cdot k \cdot n\). Its absence from the predicted spectrum is a primary open problem.

\section{Part V: \(\hbar\) as a Seam Constant}
\label{sec:org63ff2e1}

The Planck-unit route to \(\hbar\) is circular. By construction, \(E_\text{Planck} \equiv \sqrt{\hbar c^5/G}\) and \(t_\text{Planck} \equiv \sqrt{\hbar G/c^5}\), so \(E_\text{Planck} \cdot t_\text{Planck} = \hbar\) identically. Writing \(\hbar = (E_\text{Planck} \cdot t_\text{Planck})/2\pi\) yields \(\hbar/(2\pi)\), which is off by a factor of \(2\pi \approx 6.28\). This is not a derivation.

The non-circular route uses three classically measurable quantities --- proton rest mass \(m_p\), proton charge radius \(r_p\), and \(c\) --- combined with the triangle node count \(n = 3\):

\begin{equation}
\hbar = \frac{m_p \cdot c \cdot k(\gamma) \cdot r_p}{6\pi}
\end{equation}

Numerically, with \(k \approx 4.67\): \(\hbar_\text{derived} = 1.057 \times 10^{-34}\) J\(\cdot\)s vs.$\backslash$ measured \(1.055 \times 10^{-34}\) J\(\cdot\)s (0.9$\backslash$% agreement). The geometric factor \(k(\gamma)\) is not yet derived from \(\gamma\) analytically; it is fitted from the proton size. The formula is therefore an approximate consistency relation, not a prediction, until \(k(\gamma)\) is derived independently.

\section{Part VI: Cosmological Evolution and the Neutron}
\label{sec:orgab92c29}

\subsection{Cosmological Evolution of the Proton Mass}
\label{sec:orgfb6cc7e}

As the universe expands, \(R_H\) increases, \(E_\text{Hubble}\) decreases, and \(\varepsilon\) decreases:

\begin{equation}
\varepsilon(t) = \bigl(E_\text{Planck}^2 \cdot E_\text{Hubble}(t)\bigr)^{1/3}
\end{equation}

The proton mass \(m_p = 3\varepsilon\) slowly decreases over cosmic time. Rate: \(dm_p/dt = m_p \cdot H_0/3 \approx 10^{-29}\) kg/s. Unmeasurably small on human timescales; in principle detectable over cosmological timescales. This is a genuine prediction of this description.

\subsection{The Neutron Mass}
\label{sec:orgc817161}

The neutron (939.6 MeV) differs from the proton by one perturbed edge state in the triangle. In Planck units, the perturbation:

\begin{equation}
\delta = \frac{m_\text{neutron} - m_\text{proton}}{\varepsilon} = \frac{1.293\ \text{MeV}}{312.8\ \text{MeV}} = 0.00413
\end{equation}

is the fractional asymmetry of one edge state in iota-scale units. Deriving this value from the difference in up and down quark negotiation states is an open calculation.

\section{Summary}
\label{sec:org3db995d}

\begin{defbox}
\textbf{Two inputs:} iota size $1 =$ Planck length; Hubble scale $R_H$. \\[4pt]
\textbf{From the triangle topology alone, without postulation:}
\begin{itemize}\setlength\itemsep{2pt}
\item Color charge --- three distinguishable node positions in the closed loop
\item Spin-$\frac{1}{2}$ --- double degeneracy of $\lambda = 3$ mode of the triangle Laplacian
\item Confinement --- modes cannot exist outside their defining graph
\item $\varepsilon = 312.8$ MeV --- from proton mass; confirmed within 10.3\% by geometric mean formula
\item 2:1 weighting --- two iota participants per binary negotiation, one universal boundary condition
\item Hadron spectrum --- $E(n,k) = \varepsilon \cdot k \cdot n$ for complete graphs
\item Proton mass evolution --- $\varepsilon(t)$ decreases as the universe expands
\end{itemize}
\end{defbox}

\section{Document Series and Cross-References}
\label{sec:orgaefad06}

This is Document 2 of 5 in The Medium series.

\begin{center}
\begin{tabular}{clp{9cm}}
\toprule
\textbf{Doc} & \textbf{File} & \textbf{Content}\\[0pt]
\midrule
0 & medium\_doc0\_introduction.org & Personal introduction; the question behind the description\\[0pt]
1 & medium\_doc1\_definition.org & The Medium: Foundation Definition (primitives, three levels)\\[0pt]
2 & medium\_doc2\_unified.org & \emph{This paper}\\[0pt]
3 & medium\_doc3\_lagrangian.org & Lagrangian of the Negotiation Field; QFT connection; SU(3)\\[0pt]
4 & medium\_doc4\_hamiltonian.org & Hamiltonian, energy spectrum, canonical quantization, ADM\\[0pt]
5 & medium\_doc5\_G\_H0.org & Deriving \(G\) and \(H_0\) from substrate primitives\\[0pt]
\bottomrule
\end{tabular}
\end{center}

All documents: Zenodo DOI: 10.5281/zenodo.18927154 \(\cdot\) ORCID: 0009-0001-0328-2299 \(\cdot\) github.com/bkdepuy/The-Medium

\section{License}
\label{sec:org00ca507}

CC BY 4.0. You are free to share and adapt this material for any purpose, provided you give appropriate credit to the original author.
\end{document}
